\documentclass[../../main.tex]{subfiles}% 注意这里的文件路径不能用 ./main.tex ,否则用latexmk编译子文件会报错
\graphicspath{{\subfix{./image/}}} % 指定图片目录,后续可以直接使用图片文件名
% 注意这里的文件路径不能用 ../../image/ ,否则用latexmk编译子文件会报错

% 例如:
% \begin{figure}[H]
% \centering
% \includegraphics[scale=0.3]{图.png}
% \caption{}
% \label{figure:图}
% \end{figure}
% 注意:上述\label{}一定要放在\caption{}之后,否则引用图片序号会只会显示??.

\begin{document}


\section{第二基本形式}
设$\bm{r}(u,v)$是正则参数曲面，单位法向量$\bm{n}(u,v)$，则$\mathrm{d}\bm{n}=-\bm{r}_u\mathrm{d}u-\bm{r}_v\mathrm{d}v$，定义第二基本形式的系数
\[
L=-\bm{n}_u\cdot\bm{r}_u=\bm{n}\cdot\bm{r}_{uu},\quad M=-\bm{n}_u\cdot\bm{r}_v=\bm{n}\cdot\bm{r}_{uv},\quad N=-\bm{n}_v\cdot\bm{r}_v=\bm{n}\cdot\bm{r}_{vv}.
\]

\begin{definition}
二次微分形式
\[
\mathrm{II}=L(\mathrm{d}u)^2+2M\mathrm{d}u\mathrm{d}v+N(\mathrm{d}v)^2
\]
称为曲面的**第二基本形式**。
\end{definition}

第二基本形式的几何意义是曲面在点$p$的切平面与邻近点的有向距离的两倍，即
\[
\delta=-\dfrac{1}{2}\mathrm{II}(\mathrm{d}u,\mathrm{d}v).
\]

\begin{example}
球面$\bm{r}(u,v)=(a\cos u\cos v,a\cos u\sin v,a\sin u)$的第二基本形式为
\[
\mathrm{II}=-a\left((\mathrm{d}u)^2+\cos^2u(\mathrm{d}v)^2\right).
\]
\end{example}

\section{法曲率}
\begin{definition}
曲面在点$p$沿切方向$(\mathrm{d}u,\mathrm{d}v)$的**法曲率**定义为
\[
\kappa_n=\dfrac{\mathrm{II}(\mathrm{d}u,\mathrm{d}v)}{\mathrm{I}(\mathrm{d}u,\mathrm{d}v)}.
\]
\end{definition}

法曲率的符号表示曲面的弯曲方向：$\kappa_n>0$向$\bm{n}$同侧弯曲，$\kappa_n<0$向另一侧弯曲，$\kappa_n=0$为渐近方向。

\begin{theorem}
曲面在点$p$的法曲率仅与切方向有关，与曲面上经过该点的曲线选取无关。
\end{theorem}

\begin{proof}
设曲线$\bm{r}(u(t),v(t))$，曲率$\kappa$，主法向量$\bm{\beta}$，则$\kappa_n=\kappa\cos\angle(\bm{\beta},\bm{n})$，仅依赖切方向。

\end{proof}

\section{Weingarten映射和主曲率}
\begin{definition}
切空间上的线性变换$\mathscr{W}(\mathrm{d}\bm{r})=-\mathrm{d}\bm{n}$称为**Weingarten映射**，其特征值称为曲面的**主曲率**，对应的特征方向称为**主方向**。
\end{definition}

主曲率$\kappa_1,\kappa_2$是方程
\[
\begin{vmatrix}
L-\kappa E & M-\kappa F\\
M-\kappa F & N-\kappa G
\end{vmatrix}=0
\]
的根，满足
\[
\kappa_1+\kappa_2=\dfrac{EN-2FM+GL}{EG-F^2},\quad \kappa_1\kappa_2=\dfrac{LN-M^2}{EG-F^2}.
\]

\begin{definition}
主曲率的平均值称为**平均曲率**$H=\dfrac{\kappa_1+\kappa_2}{2}$，乘积称为**高斯曲率**$K=\kappa_1\kappa_2$。
\end{definition}

\begin{theorem}
主方向彼此正交的充要条件是曲面在该点是脐点（$\kappa_1=\kappa_2$）或$F=M=0$。
\end{theorem}

\section{渐近曲线和曲率线}
\begin{definition}
满足$\mathrm{II}(\mathrm{d}u,\mathrm{d}v)=0$的切方向称为**渐近方向**，以渐近方向为切方向的曲线称为**渐近曲线**。
\end{definition}

\begin{definition}
以主方向为切方向的曲线称为**曲率线**，曲率线的微分方程为
\[
\begin{vmatrix}
(\mathrm{d}u)^2 & \mathrm{d}u\mathrm{d}v & (\mathrm{d}v)^2\\
E & F & G\\
L & M & N
\end{vmatrix}=0.
\]
\end{definition}

\begin{theorem}
曲率线网是正交网的充要条件是曲面为极小曲面或无脐点且$F=M=0$。
\end{theorem}





\end{document}